%% abstract.tex: abstract in PT and EN  (FEUP regulations)
%% -------------------------------------------------------

\chapter*{Resumo}

A pontuação automática de alvos de papel é um requisito importante no tiro desportivo,
particularmente em contextos de treino, onde o feedback em tempo real pode apoiar o
desenvolvimento dos atletas. Embora os sistemas de pontuação eletrónica baseados em hardware
especializado ofereçam elevada precisão, o seu custo e os requisitos de infraestrutura limitam a
sua acessibilidade a clubes amadores. As abordagens baseadas em visão por computador
constituem uma alternativa de baixo custo, mas devem satisfazer os elevados requisitos de
precisão impostos pelas regras de pontuação decimal da ISSF.

Este trabalho aborda o problema da pontuação automática de alvos em papel da ISSF para as
modalidades de carabina de ar comprimido a 10\,m e pistola de ar comprimido a 10\,m, utilizando
sistemas de baixo custo baseados em câmaras. As características físicas dos alvos da ISSF e as
respetivas regras de pontuação são analisadas de forma a derivar requisitos de precisão para uma
pontuação decimal fiável. Uma revisão do estado da arte mostra que os métodos clássicos de
processamento de imagem podem atingir elevada precisão geométrica em condições controladas,
mas carecem de robustez, enquanto as abordagens baseadas em deep-learning
proporcionam deteção fiável, mas precisão de localização insuficiente quando utilizadas de forma
isolada.

Com base nesta análise, é proposta uma arquitetura de sistema híbrida que combina deteção
baseada em aprendizagem com refinamento geométrico clássico. É concebido um pipeline modular
de processamento de imagem para suportar a deteção robusta do alvo e dos impactos, o tratamento
de impactos sobrepostos e a estimativa precisa dos centros, permitindo uma pontuação em
conformidade com as regras da ISSF em tempo quase real. A abordagem proposta estabelece uma
base para o desenvolvimento e avaliação de um sistema de pontuação automática de baixo custo
para alvos em papel da ISSF durante a fase de dissertação.

\chapter*{Abstract}

Automatic scoring of paper targets is an important requirement in sport shooting, particularly
in training environments where real-time feedback can support athlete development. While
electronic scoring systems based on specialized hardware provide high accuracy, their cost and
infrastructure requirements limit their accessibility to amateur clubs. Computer-vision-based
approaches offer a low-cost alternative but must satisfy the high accuracy requirements
imposed by ISSF decimal scoring rules.

This work addresses the problem of automatic scoring of ISSF 10\,m air rifle and 10\,m air pistol
paper targets using low-cost camera-based systems. The physical characteristics of ISSF targets
and the scoring rules are analyzed to derive accuracy requirements for reliable decimal
scoring. A review of the state of the art shows that classical image-processing methods can
achieve high geometric precision under controlled conditions but lack robustness, while deep
learning approaches provide reliable detection but insufficient localization accuracy when used in
isolation.

Based on this analysis, a hybrid system architecture is proposed that combines learning-based
detection with classical geometric refinement. A modular image-processing pipeline is designed
to support robust target and bullet-hole detection, overlap handling and precise center
estimation, enabling ISSF-compliant scoring in near real time. The proposed approach establishes
a foundation for the development and evaluation of a low-cost automatic scoring system for
ISSF paper targets during the dissertation phase.

